% ===================================================================
%  نقشہ نمبر 3 — نِکپݨ تے جوانی۔
%  Traduction 2026 d'apres texto/io, controlee sur texto/fr, texto/en
%  et texto/hi. Consigne : voir texto/pnb/10-naqsha-01.tex.
%
%  Deux scenes, et l'ido subdivise beaucoup plus que le francais :
%  huit des intertitres n'ont pas de vis-a-vis francais. Le pendjabi
%  suit l'ido, comme les huit autres traductions.
% ===================================================================

%%K t03-apar-1 apar
\VUsaut{3.0mm}
\VUtitre{15.0pt}{60}{نقشہ نمبر 3}
\VUsaut{1.6mm}
\VUtitre{11.4pt}{0}{نِکپݨ تے جوانی۔}
\VUsaut{3.4mm}

%%K t03-apar-2 apar
(نقشے 3 تے 4 ایس طرحاں بݨائے گئے نیں پئی چار \VUgras{ٹبری دے منظراں} وچ
\VUgras{رُتاں}، \VUgras{عمراں}، \VUgras{ٹبر}، \VUgras{کپڑے} تے \VUgras{صحت} دی
مُڈھلی لفظاولی سامݨے آ جاوے۔)
\VUsaut{4.0mm}

%%K t03-tit-1 sub
\VUcentre{12.0pt}{0}{\textit{پہلا منظر۔}}
\VUsaut{3.0mm}

%%K t03-tit-2 sub
\VUpk{10.2pt}{40}{گراں وچ اک بپتسمہ۔}
\VUsaut{2.4mm}

%%K t03-tit-3 sub
\VUpk{10.2pt}{40}{بہار۔}
\VUsaut{3.0mm}

%%K t03-c1-01-1 p
1. --- اسیں \VUgras{بہار} وچ آں، \VUgras{مارچ}، \VUgras{اپریل} یا \VUgras{مئی}
دے مہینے وچ، \VUgras{ہفتے} دے کِسے دن --- \VUgras{پیر}، \VUgras{منگل} یا
\VUgras{بُدھ}۔ سویرے دے \VUgras{دس وجے} نیں۔

%%K t03-c1-02-1 p
2. --- \VUgras{موسم} سوہݨا اے، \VUgras{ہوا} نکھری ہوئی۔ سورج چمک رہیا اے۔ کجھ
نِکے نِکے \VUgras{بدل} \textsuperscript{(2)} نیلے \VUgras{اسمان}
\textsuperscript{(1)} وچ اُٹھ رہے نیں۔

%%K t03-c1-03-1 p
3. --- \VUgras{رُکھاں} اُتے ہالے \VUgras{پتّر} تقریباً نئیں۔ پتلے
\VUgras{سفیدے} \textsuperscript{(3)} تے اُچے \VUgras{چنار} \textsuperscript{(17)}
اُتے ہُݨ تیکر کلیاں ای نیں۔

%%K t03-c1-04-1 p
4. --- \VUgras{ابابیلاں} \textsuperscript{(4)} مُڑ آئیاں نیں، تے خوشی نال کوکدیاں
ہوئیاں اوس \VUgras{چوٹی} \textsuperscript{(7)} دے چوہاں پاسیں چکر کڈ رہیاں نیں
جیہڑی \VUgras{گھنٹہ گھر} \textsuperscript{(5)} دی اے، تے اوہ گھنٹہ گھر گراں دے
\VUgras{گرجے} \textsuperscript{(6)} اُتے کھڑا اے۔

%%K t03-tit-4 sub
\VUsaut{3.4mm}
\VUpk{10.2pt}{40}{گرجا۔}
\VUsaut{3.0mm}

%%K t03-c1-05-1 p
5. --- ایہہ گرجا بہت سادہ اے، اپݨے وڈے \VUgras{بُوہے} \textsuperscript{(13)} تے
اپݨی وڈی \VUgras{گھڑی} \textsuperscript{(8)} نال۔ \VUgras{نِکی سُوئی}
\textsuperscript{(9)} X دے ہندسے اُتے اے : اوہ \VUgras{گھنٹے} دسدی اے۔
\VUgras{وڈی والی} \textsuperscript{(10)} XII (دوپہر) اُتے اے؛ اوہ \VUgras{منٹ}
دسدی اے۔

%%K t03-c1-06-1 p
6. --- منارے وچ \VUgras{گھنٹے} \textsuperscript{(11)} خوشی نال وج رہے نیں۔ چھت دے
سرے اُتے پتھر دی اک نِکی \VUgras{صلیب} \textsuperscript{(12)} کھڑی اے۔

%%K t03-c1-07-1 p
7. --- گرجے وچ نِرا وڈا \VUgras{بُوہا} \textsuperscript{(13)} ای نئیں، بغل وچ اک
نِکا دروازہ وی اے۔ اوسے دروازے وچوں \VUgras{پادری} \textsuperscript{(15)}، اپݨے
\VUgras{چوغے} وچ، \VUgras{کپڑیاں والے کمرے} \textsuperscript{(14)} توں نکل کے
\VUgras{پادری دے گھر} \textsuperscript{(19)} ول جاندے نیں۔

%%K t03-c1-08-1 p
8. --- اوہناں کول ایہہ رہی \VUgras{دُدھ والی} \textsuperscript{(24)} اپݨے
\VUgras{کھوتے} \textsuperscript{(23)} نال۔ اوہ شہروں مُڑ رہی اے، جِتھے نِکیاں لئی
اپݨا چنگا دُدھ پہنچا آئی اے۔

%%K t03-tit-5 sub
\VUsaut{4.0mm}
\VUpk{10.2pt}{40}{پھلاں دا باغ۔}
\VUsaut{3.4mm}

%%K t03-c1-09-1 p
9. --- اُستاد الیبورون (کھوتا) ایس سویر دی دوڑ توں بہت خوش اے۔ اوہ گوانڈھ دے
\VUgras{باغ} دے اوہناں \VUgras{پھُلاں} \textsuperscript{(20)} توں مہکی ہوا آنند
نال اندر کھِچدا اے جیہڑے \VUgras{پھلاں والے رُکھاں} \textsuperscript{(18)} نوں
ڈھکے ہوئے نیں۔ اوہ \VUgras{آڑو دے رُکھ} \textsuperscript{(18)} تے
\VUgras{خوبانی دے رُکھ} \textsuperscript{(19)} بالکل گلابی تے چٹے نیں، تے چنگی
فصل دی آس بنھاندے نیں۔

%%K t03-c1-10-1 p
10. --- ایس دوران نِکیاں \VUgras{مکھیاں} \textsuperscript{(22)}، جیہڑیاں
\VUgras{چھتے} \textsuperscript{(21)} توں آندیاں نیں، اوتھے جا کے گُݨگُݨاندیاں
اپݨا مِٹھا \VUgras{شہد} کٹھا کردیاں نیں۔

%%K t03-c1-11-1 p
11. --- پر ایہہ کیہ ہو رہیا اے؟ گراں دے نِکے دوڑدے ہوئے آ رہے نیں تے ڈریاں
ہوئیاں \VUgras{بطخاں} \textsuperscript{(25)} نوں کھِلار رہے نیں۔ ایہہ اوہ جلوس اے
جیہڑا منشی دے پُتر دے \VUgras{بپتسمے} توں بعد گرجے وچوں نکل رہیا اے۔

%%K t03-c1-12-1 p
12. --- نِکا جیمز، اپݨے \VUgras{پنگھوڑے} \textsuperscript{(26)} وچ، گھنٹیاں دی
خوش آواز نال جاگ پیندا اے، تے اوہدی بھین میڈلن، جیہڑی جلوس ویکھݨا چاہندی اے،
اپݨی ماں نوں آکھدی اے پئی آ کے دہلیز اُتے ای اوہدے وال سنوار دیوے تے اوہنوں
کپڑے پوا دیوے۔

%%K t03-c1-13-1 p
13. --- اوہدی ماں من جاندی اے۔ اوہ اپݨا \VUgras{سر دا کپڑا} \textsuperscript{(28)}،
اپݨی \VUgras{چولی} \textsuperscript{(29)} تے اپݨا \VUgras{پرݨا} \textsuperscript{(30)}
پا کے دروازے دے سامݨے اک نِکی کرسی اُتے آ بہندی اے۔ میڈلن اُتے نِرا اوہدا
\VUgras{پیٹی کوٹ} \textsuperscript{(31)} تے اوہدی \VUgras{کُڑتی}
\textsuperscript{(32)} اے۔

%%K t03-c1-14-1 p
14. --- اوہ کِنّی خوش اے! اوہ اپݨے پیارے پھُل بھُل جاندی اے — اپݨے
\VUgras{لونگ دے پھُل} \textsuperscript{(34)}، جیہناں اُتے \VUgras{تِتلیاں}
\textsuperscript{(35)} بہندیاں نیں، اپݨے \VUgras{ٹیولپ} \textsuperscript{(36)}،
اپݨیاں \VUgras{وادی دیاں سوسناں} \textsuperscript{(37)} جیہڑیاں \VUgras{گملیاں}
\textsuperscript{(38)} وچ نیں، تے اوہ گملے \VUgras{کندھ} \textsuperscript{(33)} اُتے
رکھے نیں، تے اپݨے \VUgras{گلاب} \textsuperscript{(39)} جیہڑے اِنّی سوہݨی واڑ
بݨاندے نیں۔ اوہ جلوس دیاں بیبیاں دے قیمتی کپڑے تک رہی اے۔

%%K t03-c1-15-1 p
15. --- گراں دے مُنڈے کپڑیاں ول گھٹ دھیان دیندے نیں۔ اوہ \VUgras{مِٹھیاں}
\textsuperscript{(55)} یا \VUgras{سکے} \textsuperscript{(52)} لبھݨ لئی اگے لپکدے تے
دھکم دھکا کردے نیں۔ اوہناں وچوں اک رو \textsuperscript{(40)} رہیا اے۔

%%K t03-c1-16-1 p
16. --- دوجے نے \VUgras{کُتے} \textsuperscript{(42)} دا پنجہ مِدھ دِتا اے، جیہڑا
کوں کوں کردا نس جاندا اے تے \VUgras{کُکڑی} \textsuperscript{(43)} تے اوہدے
\VUgras{چوچیاں} \textsuperscript{(44)} نوں اُڈا دیندا اے۔

%%K t03-tit-6 sub
\VUsaut{5.0mm}
\VUpk{10.2pt}{40}{بپتسمے دا جلوس۔}
\VUsaut{4.0mm}

%%K t03-c1-17-1 p
17. --- جلوس دے اگے اگے \VUgras{دائی} \textsuperscript{(45)} ٹُر رہی اے۔ اوہدے سر
اُتے لمیاں فیتیاں والی سوہݨی \VUgras{ٹوپی} \textsuperscript{(46)} اے۔ اوہ
\VUgras{نویں جمے بال} \textsuperscript{(47)} نوں چُکی ہوئی اے، جیہڑا سارے دا سارا
\VUgras{لیس} \textsuperscript{(48)} وچ لپیٹیا ہویا اے۔

%%K t03-c1-18-1 p
18. --- دائی دے پچھے \VUgras{دھرم پیو} \textsuperscript{(49)} تے
\VUgras{دھرم ماں} \textsuperscript{(53)} آندے نیں۔ اوہ مِٹھیاں تے سکے ونڈ رہے نیں۔
دھرم پیو دے ہتھ وچ \VUgras{دستانے} \textsuperscript{(51)} تے سر اُتے
\VUgras{اُچی ہیٹ} \textsuperscript{(50)} اے۔ دھرم ماں ہتھ وچ اک سوہݨا
\VUgras{پھُلاں دا گُلدستہ} \textsuperscript{(54)} چُکی ہوئی اے۔

%%K t03-c1-19-1 p
19. --- اوہناں دے پچھے سانوں بال دے \VUgras{چاچے} \textsuperscript{(56)} تے
\VUgras{چاچی} \textsuperscript{(58)} وکھائی دیندے نیں۔ چاچی دے سر اُتے سجیلی
\VUgras{ہیٹ} \textsuperscript{(59)} اے، چاچے اُتے کالا \VUgras{لما کوٹ}
\textsuperscript{(57)} تے چٹی واسکٹ۔ فیر \VUgras{چچیرا بھرا} \textsuperscript{(60)}
تے اوہدی \VUgras{چچیری بھین} \textsuperscript{(61)} آندے نیں۔ اوہ نویں جمے بال دی
\VUgras{بھین} \textsuperscript{(62)}، نِکی برتھا، دا ہتھ پھڑی ہوئی اے۔

%%K t03-c1-20-1 p
20. --- برتھا دا دوجا \VUgras{بھرا} \textsuperscript{(63)} پچھے ٹُر رہیا اے۔ اوہ
اک \VUgras{رشتے دار} \textsuperscript{(64)} نوں ہتھ دِتا ہویا اے۔ ٹبر دے
\VUgras{یار} \textsuperscript{(65)} اوہناں توں بعد آندے نیں۔

%%K t03-tit-7 sub
\VUsaut{4.6mm}
\VUpk{10.2pt}{40}{تماشبین۔}
\VUsaut{3.4mm}

%%K t03-c1-21-1 p
21. --- جلوس کول ایہہ رہیا اک \VUgras{منگتا} \textsuperscript{(66)}، اپݨی
\VUgras{بیساکھی} \textsuperscript{(67)} دے سہارے۔ اوہ خیرات منگ رہیا اے۔

%%K t03-c1-22-1 p
22. --- دوجے پاسے اک بہت \VUgras{ادب والا} \textsuperscript{(68)} مُنڈا اے۔ اوہ
اپݨی ٹوپی چُک کے جاݨ پچھاݨ والیاں نوں \VUgras{« سلام »} آکھدا اے۔

%%K t03-c1-23-1 p
23. --- گراں والے بپتسمے دا جلوس نکلدا ویکھݨ اپݨے گھراں توں باہر آ گئے نیں۔
سانوں اپݨی \VUgras{سُوتی ٹوپی} پائی اک \VUgras{گرانویں بندہ} \textsuperscript{(69)}
وکھائی دیندا اے تے اوہدے کول اک \VUgras{گرانویں زنانی} \textsuperscript{(70)}۔ فیر
اک \VUgras{بُڈھی} \textsuperscript{(71)}، جیہڑی اپݨی \VUgras{سوٹی}
\textsuperscript{(72)} دے سہارے اے۔ چھیکر \VUgras{چکی والا} \textsuperscript{(73)}
اپݨی چوڑی \VUgras{نمدے دی ہیٹ} \textsuperscript{(74)} وچ، تے \VUgras{کھڑاواں}
\textsuperscript{(75)} پائی اک جوان گرانویں زنانی، جیہڑی اپݨے بال نوں ٹُرنا سکھا
رہی اے۔

%%K t03-c1-24-1 p
24. --- منظر بہت جیوندا جاگدا اے۔
\VUsaut{4.0mm}

%%K t03-tit-8 sub
\VUcentre{12.0pt}{0}{\textit{دوجا منظر۔}}
\VUsaut{3.0mm}

%%K t03-tit-9 sub
\VUpk{10.2pt}{40}{عوامی میلہ۔}
\VUsaut{2.4mm}

%%K t03-tit-10 sub
\VUpk{10.2pt}{40}{پاݨی دیاں کھیڈاں تے بیڑیاں دی دوڑ۔}
\VUsaut{3.0mm}

%%K t03-c2-01-1 p
1. --- \VUgras{گرمی} آ گئی اے۔ \VUgras{جون} (جولائی یا \VUgras{اگست}) دے مہینے
دا تپدا سورج جواناں نوں نہاؤݨ تے بیڑی چلاؤݨ لئی للچاندا اے۔

%%K t03-c2-02-1 p
2. --- دریا دے سجے کنڈھے اُتے ملاح پاݨی دیاں کھیڈاں تے بیڑیاں دی دوڑ وچ اُترن
لئی کپڑے لاہ رہے نیں۔

%%K t03-c2-03-1 p
3. --- اوہناں وچوں اک اپݨے \VUgras{جُتے} \textsuperscript{(1)} لاہ رہیا اے؛ اوہ
اوہناں دے تسمے (بٹن) کھول رہیا اے۔ اوہدے دو ساتھی پہلاں ای تیار نیں۔ ہُݨ اوہناں
اُتے نِری اپݨی \VUgras{بنیان} \textsuperscript{(2)} تے \VUgras{تیرن دا جانگیا}
\textsuperscript{(3)} اے۔ اوہ اپݨے ساتھیاں دی اُڈیک کردے گلاں کر رہے نیں۔ اوہ
دوجے کنڈھے اُتے لگے میلے تے جشن دی گل کر رہے نیں۔

%%K t03-c2-04-1 p
4. --- اوہناں کول تھلے اوہناں دے ساتھیاں دے \VUgras{کپڑے} پئے نیں : اک
\VUgras{جوڑی} \VUgras{بوٹ} \textsuperscript{(4)}، \VUgras{جُتے}
\textsuperscript{(5)}، \VUgras{جُرابا} \textsuperscript{(6)}، \VUgras{نمدے}
\textsuperscript{(7)} تے \VUgras{تیلیاں} \textsuperscript{(8)} دیاں ہیٹاں تے اک
\VUgras{ٹائی} \textsuperscript{(9)}۔

%%K t03-c2-05-1 p
5. --- اک جوان نے اپݨا \VUgras{کوٹ} \textsuperscript{(10)} ہوشیاری نال تہہ کیتا
اے۔ اوہ اوہنوں اک تولیے اُتے رکھ رہیا اے، جیدے وچ اوہدا گوانڈھی تھوڑی دیر مگروں
اوہناں دوہاں دے \VUgras{سوٹ} بنھ دیوے گا۔

%%K t03-c2-06-1 p
6. --- چھیواں مُنڈا پچھے رہ گیا اے۔ اوہدے سر اُتے ہالے وی اوہدی \VUgras{ٹوپی}
\textsuperscript{(11)} اے۔ اوہنے نہ اپݨی \VUgras{قمیض} \textsuperscript{(12)} لاہی
اے، نہ اپݨا \VUgras{کالر} \textsuperscript{(13)}، نہ اپݨے \VUgras{کف}
\textsuperscript{(14)}۔ اوہ ہتھ وچ اپݨی \VUgras{واسکٹ} \textsuperscript{(17)} چُکی
ہوئی اے تے ہالے وی اپݨی \VUgras{پتلون} \textsuperscript{(16)} تے اپݨیاں
\VUgras{گیلساں} \textsuperscript{(15)} پائی ہوئی اے۔ اوہ اگلی دوڑ لئی ضرور تیار نہ
ہو سکے گا۔

%%K t03-c2-07-1 p
7. --- جواناں کول \VUgras{اُٹ کٹاریاں} \textsuperscript{(18)} نال گھِری اک پگڈنڈی
دریا دے کنڈھے اُتردی اے، جِتھے بُڈھا جیمز \VUgras{کانیاں} \textsuperscript{(19)}
وچکار شام نوں چھڈی جاݨ والی \VUgras{آتشبازی} \textsuperscript{(20)} تیار کر رہیا
اے۔

%%K t03-tit-11 sub
\VUsaut{4.4mm}
\VUpk{10.2pt}{40}{دوڑ۔}
\VUsaut{3.4mm}

%%K t03-c2-08-1 p
8. --- دریا اُتے اپݨیاں چھتریاں دی چھاں وچ کجھ بیبیاں اک \VUgras{ڈونگی}
\textsuperscript{(21)} وچ نکلیاں نیں۔ اوہ دوڑ ویکھ رہیاں نیں، جیہڑی بہت رومانچک
اے۔ ہر \VUgras{بیڑی} \textsuperscript{(22)} وچ چار \VUgras{چپو والے}
\textsuperscript{(24)} پوری تاکت نال کھِچ رہے نیں، جدوں کہ \VUgras{ملاح}
\textsuperscript{(26)} \VUgras{پتوار} \textsuperscript{(25)} پھڑی نازک بیڑی نوں سیدھ
دیندا اے۔ \VUgras{چپو} \textsuperscript{(23)} تال نال پاݨی اُتے پیندے نیں، تے
ہولیاں بیڑیاں چکرا دیݨ والی چال نال دوڑدیاں نیں۔

%%K t03-c2-09-1 p
9. --- کجھ جوان پُل دے تھمّ کول \VUgras{چِکݨے بانس} \textsuperscript{(28)} دا
انعام جِتݨ وچ لگے نیں۔ اوہناں وچوں اک لچکدے تِلکݨے بانس اُتے ہوشیاری نال سرکدا
ہویا سرے اُتے بنھی نِکی \VUgras{جھنڈی} \textsuperscript{(29)} تیکر اپڑنا چاہندا اے۔
اوہدا نِگھرا \VUgras{مخالف} \textsuperscript{(30)} پاݨی وچ ڈِگ پیا اے۔ اوہ تیر کے
کنڈھے مُڑ رہیا اے۔

%%K t03-c2-10-1 p
10. --- وڈے مُنڈے \VUgras{پاݨی دی لڑائی} \textsuperscript{(32)} کھیڈ رہے نیں،
\VUgras{گھاٹ} \textsuperscript{(33)} کول۔ اک وڈی \VUgras{بِھیڑ} \textsuperscript{(31)}
\VUgras{پُل} \textsuperscript{(27)} دی \VUgras{مُنڈیر} اُتے جھُکی تے
\VUgras{کنڈھے} اُتے ٹھسا ٹھس بھری ایہہ ساریاں کھیڈاں ویکھ رہی اے۔ پر کجھ تماشبین
مُڑ کے \VUgras{چڑھن دے بانس} \textsuperscript{(45)} دی کھیڈ ویکھدے نیں، یا
\VUgras{انعام ونڈائی} ول جاندے \VUgras{مئیر دے جلوس} نوں تکدے نیں۔

%%K t03-c2-11-1 p
11. --- سب توں پہلاں ایہہ رہے کجھ \VUgras{چھوکرے} \textsuperscript{(35)}، جیہڑے
\VUgras{ساز والیاں} \textsuperscript{(36)} دے اگے اگے دوڑ رہے نیں؛ فیر آندے نیں
\VUgras{مئیر} \textsuperscript{(37)}، اوہناں دے نائب تے نِکے شہر دی
\VUgras{شہری کونسل}۔ سب توں پچھے \VUgras{اگ بجھاؤ والے} \textsuperscript{(38)}
ٹُردے نیں۔

%%K t03-tit-12 sub
\VUsaut{5.0mm}
\VUpk{10.2pt}{40}{میلہ۔}
\VUsaut{3.6mm}

%%K t03-c2-12-1 p
12. --- \VUgras{میلے دے میدان} \textsuperscript{(39)} اُتے وڈی بِھیڑ اے۔ لوک کئی
طرحاں دے \VUgras{تنبو} ویکھݨ آندے نیں : \VUgras{لکڑ دے گھوڑے} \textsuperscript{(40)}،
\VUgras{سرکس} \textsuperscript{(41)}، جِتھے کھیڈ شروع وی ہو چُکی اے؛
\VUgras{عوامی ناچ} \textsuperscript{(42)}، جِتھے شہر دے جوان مُنڈے تے کُڑیاں
\VUgras{نچّݨ} دا سواد لے رہے نیں۔

%%K t03-c2-13-1 p
13. --- سرے اُتے سانوں \VUgras{اکھاڑا} \textsuperscript{(44)} وکھائی دیندا اے،
جِتھے سورما ہسپانوی \VUgras{بلد باز} گُصیلے \VUgras{بلد} \textsuperscript{(45)} نال
لڑدا اے؛ فیر \VUgras{جھوٹے والی ریل} \textsuperscript{(49)} تے
\VUgras{نشانہ بازی دا تنبو} \textsuperscript{(46)}، جِتھے لوک \VUgras{بندوق}
چلاؤݨ تے \VUgras{نشانے} اُتے مارن دی مشق کردے نیں۔

%%K t03-c2-14-1 p
14. --- کول ای ایہہ رہیا \VUgras{میلے دا رنگ منچ} \textsuperscript{(47)}، جِتھے
\VUgras{ٹھٹھول} تے \VUgras{ناٹک} کھیڈے جاندے نیں۔ اوہدے بالکل کول
\VUgras{دیو} \textsuperscript{(48)} تے \VUgras{بونے} دا تنبو اے۔ پہلے دا
\VUgras{قد} دو میٹر پندرہ اے؛ دوجا مشکل نال اک بال جِنّا اُچا اے۔

%%K t03-c2-15-1 p
15. --- چھیکر ایہہ رہی وڈی \VUgras{چڑیا گھر دی گڈی} \textsuperscript{(50)}، اپݨے
\VUgras{جنگلی جانوراں} نال — اپݨے \VUgras{شیر} \textsuperscript{(51)} نال، جیدے
\VUgras{پنجے} بڑے تکڑے نیں، اپݨے \VUgras{ببر شیر} \textsuperscript{(52)} نال، جیدی
\VUgras{ایال} سنی اے، اپݨیاں دو \VUgras{جرافاں} \textsuperscript{(53)} تے اپݨے
\VUgras{ہاتھی} \textsuperscript{(54)} نال، جیہڑا اپݨی \VUgras{سُونڈ} اِنّی
ہوشیاری نال کم وچ لیاندا اے۔

%%K t03-c2-16-1 p
16. --- \VUgras{جانور سِدھاؤݨ والا} \textsuperscript{(55)}، جیہڑا ایہناں جانوراں
نوں سکھاندا اے، تنبو دے دروازے اُتے کھڑا اے۔
\VUsaut{6.0mm}
\VUfilet{22mm}
